\documentclass[a4paper,10pt]{article}
\usepackage[utf8]{inputenc}
\usepackage[T1]{fontenc}
\usepackage[french]{babel}
\usepackage{amsmath,amssymb}
\usepackage{geometry}
\usepackage{array,tabularx}
\usepackage{tikz}
\geometry{margin=1.6cm}
\pagestyle{empty}
\newcommand{\ligne}{\par\vspace{0.55cm}\noindent\dotfill\par}
\newcommand{\carreaux}[1]{\begin{center}\begin{tikzpicture}\draw[step=0.5cm,gray!35,very thin] (0,0) grid (17,#1);\end{tikzpicture}\end{center}}
\newenvironment{exercice}[2]{\paragraph{Exercice #1 \hfill #2 points}\mbox{}\par}{\medskip}
\begin{document}
\begin{center}\Large\bfseries Devoir surveillé 12 - Grandeurs et mesures\end{center}
\vspace{2mm}
\begin{exercice}{1}{4}Convertir : $3,2$ hL en L, $4500$ mL en L, $0,08$ km$^2$ en m$^2$.\end{exercice}\begin{exercice}{2}{5}Une voiture parcourt $156$ km en $2$ h. Calculer sa vitesse moyenne.\end{exercice}\begin{exercice}{3}{5}Un robinet débite $18$ L par minute. Combien de temps faut-il pour remplir une cuve de $270$ L ?\end{exercice}\begin{exercice}{4}{6}Un plan est à l'échelle $1/200$. Une pièce mesure $4,8$ cm sur le plan. Quelle est sa longueur réelle ?\end{exercice}\end{document}
